\section*{Notation and proof objects}
\addcontentsline{toc}{section}{Notation and proof objects}
\label{sec:conventions-v3}

We collect here only the conventions that are used throughout several later
sections. More local notation is introduced where it first becomes necessary.
This keeps the main argument self-contained without forcing the reader to
retain a large dictionary before the proof begins.

\paragraph{Asymptotic conventions.}
All logarithms are natural. An event holds \emph{with high probability} if its
probability tends to one as $n\to\infty$. Whenever a phase parameter is
present, $o(1)$ denotes one deterministic error sequence that is uniform over
the full phase interval, including sequences approaching either endpoint.
For integers $x,r\ge0$, we use the falling factorial
\[
  (x)_r=x(x-1)\cdots(x-r+1),
  \qquad
  (x)_0=1,
\]
and set $(x)_r=0$ for $r>x$. A quotient involving falling factorials is used
only after the denominator has been shown to be nonzero.

\paragraph{Signed witnesses and profiles.}
A \emph{signed cocoloring witness} is a partition together with an $I$- or
$K$-mark on each class. An $I$-marked class must be independent and a
$K$-marked class must be complete. The marks are counting data; they do not
define a new graph invariant. In the four-size profile every class has size at
least two for sufficiently large $n$, so a realized class cannot satisfy both
requirements. Forgetting the marks therefore recovers its cocoloring without
multiplicity.

A \emph{profile} is a finite sequence $(k_s)$, where $k_s$ is the number of
classes of size $s$. It is feasible on $n$ vertices when
\[
  \sum_s k_s=k,
  \qquad
  \sum_s s k_s=n.
\]
The ordinary first moment counts partitions into independent classes. The
signed first moment counts the same partitions together with one $I/K$ mark
per class.

\paragraph{Overlap table and support graph.}
Let $(A_a)_a$ and $(B_b)_b$ be two ordered profile partitions. Their overlap
table is
\[
  r_{ab}=|A_a\cap B_b|.
\]
Its row and column sums are the class sizes of the two partitions. It may also
be viewed as the cell-count table of a bipartite configuration model: every
underlying vertex pairs one labeled row stub with one labeled column stub.

The \emph{support graph} $H(r)$ is the simple bipartite graph with
\[
  E(H(r))=\{(a,b):r_{ab}\ge2\}.
\]
A subset of its edges is \emph{even} if every vertex has even degree. Writing
$c(H)$ for the number of connected components, including isolated vertices,
we put
\[
  \beta(H)=|E(H)|-|V(H)|+c(H).
\]
Then $\beta(H)$ is the dimension of the binary cycle space, so the number of
even edge sets is $2^{\beta(H)}$. The threshold two is exact: a cell of size
zero or one contains no edge internal to both partition classes, whereas a
cell of size at least two forces their two signs to agree.

\paragraph{Three matching levels.}
The proof uses three distinct objects that should not be conflated.
\begin{enumerate}
\item The \emph{configuration matching} is the perfect matching of all labeled
row and column stubs induced by the common vertex set.
\item A \emph{physical partial matching} is a set of selected stub pairs inside
specified overlap cells.
\item A \emph{block support} is a matching between row-class slots and
column-class slots; its edges indicate the cells declared canonical and high.
\end{enumerate}
The product formula for physical cell counts relies on the third object being
a matching. Distinct selected cells then use disjoint row and column stub
families.

\paragraph{High cells, deficits, and endpoint tables.}
Let $U$ be the largest class size in the four-size profile and put
$R_0=\lfloor U/2\rfloor$. A cell is \emph{high} when its multiplicity exceeds
$R_0$. Two high cells cannot share a row or a column, so the high cells form a
canonical block support. For a selected cell $e$, let $s_e,t_e$ be its endpoint
sizes and set
\[
  m_e=\min\{s_e,t_e\},
  \qquad
  j_e=m_e-h_e.
\]
Here $m_e$ is the full-containment multiplicity, $j_e$ is the realized
multiplicity, and $h_e$ is the deficit. Highness implies $2h_e<m_e$.

The four endpoint sizes are indexed by $i,j\in\{0,1,2,3\}$. For a block support
$P$, the endpoint table $L(P)=(L_{ij})$ records how many support edges join
endpoint types $i$ and $j$. A \emph{realized} endpoint table is one attained by
an actual finite block support; it is not an arbitrary element of
$\mathbb N^{4\times4}$.

\paragraph{Separation of accounts.}
The \emph{bare high-skeleton weight} contains the incidence probability and
local rewards of the exposed high cells. It contains neither the rewards of
unexposed cells nor the residual binary cycle-space factor. Conditional on a
high skeleton, the expectation of those remaining factors is its
\emph{residual attachment}. Sections~8 and~9 preserve this separation exactly:
a factor paid in one account is never paid again in the other.

Every load-bearing estimate follows the same order. We first state the finite
identity, then apply deterministic inequalities, then take phase asymptotics,
and only at the end deduce a probability statement. This order is useful both
for mathematical auditing and for the theorem-by-theorem formalization.
